%% =========================================================
%% Expected trace error of corrected sources.
%% Source: September 24, 2026 polynomial-PEPS preprint, Theorem 5.2.
%% This chapter does not assert the circuit expansion or network conclusion.
%% =========================================================

\chapter{Expected Trace Error of Corrected Sources}\label{ch:source_error}

The Gaussian replacement of a Schmidt source is unbiased. Its centered
entries have diagonal covariance, and fresh independent slots preserve
that covariance under tensor products. Quarter-power comparisons and
finite block aggregation then bound the expected trace error for a fixed
contraction between the complete free input spaces.
Finite branch sums and insertion of physical matrix units preserve the
required trace-norm estimates.
These are the source calculations in
\cite[Theorem~5.2]{OpenAI2026PolynomialPEPS}.

\section{An unbiased product-operator source}

\begin{definition}[Sampled Schmidt-source operators]\label{def:source_product}
    \lean{QICLean.ComplexGaussian.schmidtVector}
    \lean{QICLean.ComplexGaussian.schmidtSource}
    \lean{QICLean.ComplexGaussian.sourceU}
    \lean{QICLean.ComplexGaussian.sourceV}
    \lean{QICLean.ComplexGaussian.sampledSource}
    \lean{QICLean.ComplexGaussian.sourceCorrection}
    \lean{QICLean.ComplexGaussian.schmidtSource_apply}
    \lean{QICLean.ComplexGaussian.sampledSource_apply}
    \leanok
    Let $A,C$ be finite sets with nonnegative weights $\lambda_a,\mu_c$,
    and let $k$ be a positive integer. In the paired coordinate bases, put
    $|\eta\rangle=\sum_a\sqrt{\lambda_a}|a,a\rangle$ and
    $|\widetilde\eta\rangle=\sum_c\sqrt{\mu_c}|c,c\rangle$.
    For independent standard circular complex Gaussian coordinates
    $g^{(j)}_{ac}$, with $\mathbb E|g^{(j)}_{ac}|^2=1$, set
    \begin{align}
        (U_j)_{ac} & =(\lambda_a\mu_c)^{1/4}g^{(j)}_{ac},
        \label{eq:source_u}                                          \\
        (V_j)_{bd} & =(\lambda_b\mu_d)^{1/4}\overline{g^{(j)}_{bd}},
        \label{eq:source_v}                                          \\
        \widehat X & =\frac1k\sum_{j=0}^{k-1}U_j\otimes V_j,
        \label{eq:source_sampled}                                    \\
        X          & =|\eta\rangle\langle\widetilde\eta|,
        \label{eq:source_target}                                     \\
        \Delta     & =\widehat X-X.
        \label{eq:source_correction}
    \end{align}
    Thus $X_{(a,b),(c,d)}=\sqrt{\lambda_a\mu_c}\delta_{ab}\delta_{cd}$,
    and
    \begin{align}
        \widehat X_{(a,b),(c,d)}
         & =(\lambda_a\lambda_b\mu_c\mu_d)^{1/4}
        \frac1k\sum_{j=0}^{k-1}g^{(j)}_{ac}\overline{g^{(j)}_{bd}}.
        \label{eq:source_sampled_entry}
    \end{align}
    This is \cite[Equation~(5.6)]{OpenAI2026PolynomialPEPS}.
    % Source: eq:compression-random-source, 04-compression.tex:279--309.
    % The matrices have shape (A x A) by (C x C); the two supports may differ.
\end{definition}

\begin{theorem}[Unbiasedness and actual correction-entry covariance]
    \label{thm:source_unbiased}
    \lean{QICLean.ComplexGaussian.weighted_delta_eq_schmidtSource}
    \lean{QICLean.ComplexGaussian.sampleAverage_eq_average_sub_delta}
    \lean{QICLean.ComplexGaussian.sourceCorrection_apply}
    \lean{QICLean.ComplexGaussian.integrable_sourceCorrection_entry}
    \lean{QICLean.ComplexGaussian.integrable_sourceCorrection_entry_mul_conj}
    \lean{QICLean.ComplexGaussian.integral_sourceCorrection_entry_mul_conj}
    \lean{QICLean.ComplexGaussian.integrable_sampledSource_entry}
    \lean{QICLean.ComplexGaussian.integral_sampledSource_entry}
    \lean{QICLean.ComplexGaussian.integrable_sampledSource}
    \lean{QICLean.ComplexGaussian.integrable_sourceCorrection}
    \lean{QICLean.ComplexGaussian.integral_sampledSource}
    \lean{QICLean.ComplexGaussian.integral_sourceCorrection}
    \leanok
    \uses{def:source_product}
    For the Gaussian probability law of Definition~\ref{def:source_product},
    $\widehat X$ and $\Delta$ are integrable, with
    $\mathbb E\widehat X=X$ and $\mathbb E\Delta=0$.
    The correction entries are precisely
    \begin{align}
        \Delta_{(a,b),(c,d)}
         & =(\lambda_a\lambda_b\mu_c\mu_d)^{1/4}
        \left(\frac1k\sum_{j=0}^{k-1}g^{(j)}_{ac}
        \overline{g^{(j)}_{bd}}-\delta_{ab}\delta_{cd}\right).
        \label{eq:source_correction_entry}
    \end{align}
    Every conjugate pair product is integrable. For $p,r\in A\times A$
    and $q,t\in C\times C$,
    \begin{align}
        \mathbb E[\Delta_{pq}\overline{\Delta_{rt}}]
         & =\frac1k\sqrt{\lambda_{p_1}\lambda_{p_2}\mu_{q_1}\mu_{q_2}}
        \delta_{pr}\delta_{qt}.
        \label{eq:source_entry_covariance}
    \end{align}
    No normalization of the weight lists is required for these identities.
    These are the unbiasedness and covariance calculations in
    \cite[Equation~(5.6) and Equations~(5.7)--(5.8)]{OpenAI2026PolynomialPEPS}.
    % Source: eq:compression-random-source, eq:compression-gaussian-covariance,
    % and eq:compression-slot-variance, 04-compression.tex:279--340.
\end{theorem}

\begin{proof}\leanok
    Expanding the Kronecker product gives (\ref{eq:source_sampled_entry}).
    The quarter weight times $\delta_{ab}\delta_{cd}$ equals $X$ entrywise,
    giving (\ref{eq:source_correction_entry}). The average of the centered
    Gaussian products has mean zero and covariance $k^{-1}$ times the
    product of coordinate deltas, proving (\ref{eq:source_entry_covariance}).
    Its $L^2$ bound gives entry and pair-product integrability. Finite
    coordinate evaluation gives matrix integrability and permits the
    entrywise expectations to be combined into the matrix identities.
\end{proof}

\section{Finite branch sums and physical matrix entries}

\begin{theorem}[Complex coefficients in finite trace-norm sums]
    \label{thm:source_branch_sum}
    \lean{Matrix.rectangularTraceNorm_zero}
    \lean{Matrix.rectangularTraceNorm_add_le}
    \lean{Matrix.rectangularTraceNorm_smul}
    \lean{Matrix.rectangularTraceNorm_sum_le}
    \lean{Matrix.rectangularTraceNorm_sum_smul_le}
    \lean{Matrix.rectangularTraceNorm_conjTranspose_le}
    \lean{Matrix.rectangularTraceNorm_conjTranspose}
    \leanok
    For arbitrary rectangular complex matrices of a fixed size, the
    trace norm satisfies $\lVert0\rVert_1=0$,
    $\lVert A+B\rVert_1\leq\lVert A\rVert_1+\lVert B\rVert_1$,
    $\lVert cA\rVert_1=|c|\lVert A\rVert_1$, and
    $\lVert A^\dagger\rVert_1=\lVert A\rVert_1$.
    For a finite set $B$ of branches and complex coefficients $c_\beta$,
    \begin{align}
        \Bigl\lVert\sum_{\beta\in B}c_\beta A_\beta\Bigr\rVert_1
         & \leq\sum_{\beta\in B}|c_\beta|\lVert A_\beta\rVert_1.
        \label{eq:source_branch_sum}
    \end{align}
    Taking all $c_\beta=1$ gives the unweighted finite-sum estimate; the
    empty branch family gives zero.
    These are the trace-norm properties used in the summation leading to
    \cite[Equation~(5.19)]{OpenAI2026PolynomialPEPS}.
    % Source: eq:compression-total-error, 04-compression.tex:540--551.
\end{theorem}

\begin{proof}\leanok
    \uses{thm:random_trace_duality}
    For a contraction $V$, expand the trace pairing of the finite sum
    into $\sum_\beta\overline{c_\beta}\tr(A_\beta^\dagger V)$.
    The triangle inequality and trace-norm duality give
    (\ref{eq:source_branch_sum}), including the additive bound and zero
    case. An attaining contraction for $A$ gives the reverse scalar
    inequality, hence exact homogeneity. Adjointing the test contraction
    and cyclically permuting the trace gives
    $\lVert A^\dagger\rVert_1\leq\lVert A\rVert_1$; applying this to the
    adjoint gives equality.
\end{proof}

\begin{theorem}[Exact physical matrix-unit reinsertion]
    \label{thm:source_physical_reinsertion}
    \lean{Matrix.l2_opNorm_le_one_of_conjTranspose_mul_self_eq_one}
    \lean{Matrix.rectangularTraceNorm_isometry_sandwich}
    \lean{Matrix.basisRegisterInjection}
    \lean{Matrix.basisRegisterInjection_conjTranspose_mul_self}
    \lean{Matrix.basisRegisterInjection_sandwich_eq_single_kronecker}
    \lean{Matrix.rectangularTraceNorm_single_kronecker}
    \leanok
    Let $A$ be any rectangular complex matrix, and let $K,L$ be
    isometric embeddings of its ket and bra coordinate spaces. Thus
    $K^\dagger K=\Id$ and $L^\dagger L=\Id$; both embeddings are operator
    contractions, and
    \begin{align}
        \lVert KAL^\dagger\rVert_1 & =\lVert A\rVert_1.
        \label{eq:source_isometric_insertion}
    \end{align}
    In particular, for finite physical ket and bra registers with
    coordinates $x,y$, let $J_xv=|x\rangle\otimes v$ and
    $J_yw=|y\rangle\otimes w$. These maps are isometries and satisfy
    \begin{align}
        J_xAJ_y^\dagger                            & =|x\rangle\langle y|\otimes A,
        \label{eq:source_matrix_unit}                                               \\
        \lVert|x\rangle\langle y|\otimes A\rVert_1 & =\lVert A\rVert_1.
        \label{eq:source_physical_reinsertion}
    \end{align}
    The physical registers may differ, and no positivity or Hermiticity
    assumption on $A$ is required.
    This is the reinsertion step following
    \cite[Equation~(5.18)]{OpenAI2026PolynomialPEPS}.
    % Source: eq:compression-one-choice, 04-compression.tex:535--538.
\end{theorem}

\begin{proof}\leanok
    \uses{thm:random_exterior_contraction}
    The isometry equations imply the operator contraction bounds.
    Applying left and right trace-norm contraction gives one direction of
    (\ref{eq:source_isometric_insertion}). For the reverse direction,
    apply the adjoint contractions to $KAL^\dagger$ and use
    $K^\dagger(KAL^\dagger)L=A$.
    The identities $J_x^\dagger J_x=\Id$ and
    (\ref{eq:source_matrix_unit}) follow by coordinate evaluation.
    Substitute these isometries into (\ref{eq:source_isometric_insertion})
    to obtain (\ref{eq:source_physical_reinsertion}).
\end{proof}

\section{Covariance and weighted block errors}

\begin{theorem}[Expected trace error from diagonal covariance]
    \label{thm:source_covariance_error}
    \lean{ProbabilityTheory.sourceBlock}
    \lean{ProbabilityTheory.weightedSourceError}
    \lean{ProbabilityTheory.weightedSourceQuarter}
    \lean{ProbabilityTheory.weightedSourceError_eq_halfWeighted}
    \lean{ProbabilityTheory.quarterWeighted_source_sum_eq_transpose}
    \lean{ProbabilityTheory.rectangularTraceNorm_weightedSourceError_le}
    \lean{ProbabilityTheory.sourceProductProbability}
    \lean{ProbabilityTheory.sum_sourceProductProbability}
    \lean{ProbabilityTheory.weighted_source_block_aggregation}
    \lean{ProbabilityTheory.weighted_source_block_sum_le_one}
    \lean{ProbabilityTheory.aestronglyMeasurable_rectangularTraceNorm_sum}
    \lean{ProbabilityTheory.integrable_weightedSourceQuarter_norm}
    \lean{ProbabilityTheory.integrable_rectangularTraceNorm_weightedSourceError}
    \lean{ProbabilityTheory.integral_frobeniusNormSq_weightedSourceQuarter_eq}
    \lean{ProbabilityTheory.integral_rectangularTraceNorm_weightedSourceError_le}
    \leanok
    Let $I,L,T,U$ be finite index sets, let $p_i,\widetilde p_l$ be
    nonnegative probability lists, and let $\tau,\widetilde\tau$ be
    nonnegative diagonal trace-one matrices on the crossing-support spaces
    $T,U$. Let $O:\C^{I\times T}\to\C^{L\times U}$ be a contraction,
    with blocks $O_{li}:\C^T\to\C^U$.
    On a probability space, let the complex coefficients $z_{il}$ have
    integrable conjugate pair products and suppose, for $\kappa\geq0$,
    \begin{align}
        \mathbb E[z_{il}\overline{z_{i'l'}}]
         & =\kappa\sqrt{p_i\widetilde p_l}\delta_{ii'}\delta_{ll'}.
        \label{eq:source_covariance}
    \end{align}
    Define
    \begin{align}
        Z & =\sum_{i,l}z_{il}\tau^{1/2}O_{li}^{T}\widetilde\tau^{1/2},
        \label{eq:source_exterior_input}                               \\
        H & =\sum_{i,l}z_{il}\widetilde\tau^{1/4}O_{li}\tau^{1/4}.
        \label{eq:source_quarter_sum}
    \end{align}
    The transpose is ordinary full transpose, preserving the coefficients.
    The functions $\lVert Z\rVert_1$ and $\lVert H\rVert_2$ are integrable,
    and
    \begin{align}
        \mathbb E\lVert H\rVert_2^2
                                  & =\kappa\sum_{i,l}\sqrt{p_i\widetilde p_l}
        \lVert\widetilde\tau^{1/4}O_{li}\tau^{1/4}\rVert_2^2,
        \label{eq:source_block_moment}                                        \\
        \mathbb E\lVert Z\rVert_1 & \leq\sqrt\kappa.
        \label{eq:source_expected_error}
    \end{align}
    These are the covariance implication and block calculation in
    \cite[Equations~(5.14)--(5.18)]{OpenAI2026PolynomialPEPS}.
    % Source: eq:compression-exterior-input, eq:compression-block-second-moment,
    % eq:compression-dimension-free, eq:compression-one-choice, 04-compression.tex:480--533.
    % The covariance and pair-product integrability are explicit premises here.
\end{theorem}

\begin{proof}\leanok
    \uses{thm:random_quarter_weights, thm:random_matrix_moments}
    Quarter-power comparison and ordinary transpose give
    $\lVert Z\rVert_1\leq\lVert H\rVert_2$.
    Diagonal covariance and Theorem~\ref{thm:random_matrix_moments} give
    (\ref{eq:source_block_moment}). Put
    $R=\diag(p)\otimes\tau$ and
    $\widetilde R=\diag(\widetilde p)\otimes\widetilde\tau$.
    Their diagonal lists $p_i\tau_t$ and $\widetilde p_l\widetilde\tau_u$
    sum to one. Entrywise expansion gives the exact block aggregation
    \begin{align}
        \sum_{i,l}\sqrt{p_i\widetilde p_l}
        \lVert\widetilde\tau^{1/4}O_{li}\tau^{1/4}\rVert_2^2
         & =\lVert\widetilde R^{1/4}OR^{1/4}\rVert_2^2\leq1.
        \label{eq:source_aggregation}
    \end{align}
    Pair-product integrability makes the Gram matrix $Z^\dagger Z$
    almost everywhere measurable; taking its positive square root and
    trace gives almost everywhere strong measurability of
    $\lVert Z\rVert_1$.
    The integrable square of $\lVert H\rVert_2$ gives its integrability,
    and the pointwise comparison gives integrability of $\lVert Z\rVert_1$.
    Finally, Cauchy--Schwarz on the probability space yields
    $\mathbb E\lVert Z\rVert_1\leq\mathbb E\lVert H\rVert_2
        \leq\sqrt{\mathbb E\lVert H\rVert_2^2}\leq\sqrt\kappa$.
\end{proof}

\begin{theorem}[Trace error under the constructed independent Gaussian law]
    \label{thm:source_gaussian_error}
    \lean{QICLean.ComplexGaussian.integrable_rectangularTraceNorm_gaussianSourceError}
    \lean{QICLean.ComplexGaussian.integral_frobeniusNormSq_gaussianSourceQuarter_eq}
    \lean{QICLean.ComplexGaussian.integral_rectangularTraceNorm_gaussianSourceError_le}
    \leanok
    Let $S$ be a finite set of corrected slots, with $s=|S|$. At each slot,
    choose finite ket and bra supports $A_s,C_s$ and normalized nonnegative
    endpoint weights. The corrected-register index sets are
    $I=\prod_s(A_s\times A_s)$ and $L=\prod_s(C_s\times C_s)$.
    Construct the Gaussian coefficients and their independent
    product law as in Theorem~\ref{thm:random_gaussian_slots}.
    Let $\tau,\widetilde\tau$ be nonnegative diagonal trace-one matrices,
    and let $O$ be a contraction between the complete free registers.
    For every positive sample count $k$, define $Z,H$ by
    (\ref{eq:source_exterior_input}) and (\ref{eq:source_quarter_sum}).
    Then $\lVert Z\rVert_1$ is integrable, the second moment of $H$ is
    (\ref{eq:source_block_moment}) with $\kappa=k^{-s}$, and
    \begin{align}
        \mathbb E\lVert Z\rVert_1 & \leq k^{-s/2}.
        \label{eq:source_gaussian_error}
    \end{align}
    Here the expectation uses the constructed independent Gaussian law.
    This is \cite[Equation~(5.18)]{OpenAI2026PolynomialPEPS}.
    % Source: eq:compression-product-covariance, eq:compression-block-second-moment,
    % eq:compression-dimension-free, eq:compression-one-choice, 04-compression.tex:394--538.
    % No covariance, moment, or coefficient-measurability premise is added.
    % The circuit must still supply the fixed full-register contraction O.
\end{theorem}

\begin{proof}\leanok
    \uses{thm:random_gaussian_slots, thm:source_covariance_error}
    The actual independent Gaussian law supplies the diagonal covariance
    (\ref{eq:source_covariance}) with $\kappa=k^{-s}$ and supplies pair-product
    integrability. Normalization of the endpoint lists gives normalization
    of the product lists $p,\widetilde p$.
    Apply Theorem~\ref{thm:source_covariance_error} and use
    $\sqrt{k^{-s}}=k^{-s/2}$ for $k>0$.
\end{proof}
